% !PRINTED 706
\weberpage{706}
In order now to transform this expression for the function $\varPhi$ further
and to prepare it for computation, one has to use the theorems
on \gesperrt{Abel}ian groups and group characters
which we became acquainted with in sections 14 to 16 of this
volume.

The totality of the numbers $n$ not divisible by $q$
forms, taken with respect to the modulus $m$, an \gesperrt{Abel}ian group $\mathfrak{N}$
of degree $\mu$, whose characters $\chi(n)$ are determined in \S.\,16. There
was in this a small difference, according as $q$ is odd or
$q=2$.

1) If $q$ is odd, there is a primitive root $c$ of $m$,
and for every number $n$ there is an index $\gamma$ such that
\weq{5}{n \equiv c^{\gamma}\ (\mathrm{mod}\ m).}
If then $\Theta$ is a $\mu$th root of unity, then
\weq{6}{\chi(n) = \Theta^{\gamma},}
and one obtains all $\mu$ characters if for $\Theta$
one puts the different $\mu$th roots of unity.

If $n$ belongs to the exponent $f$, then $e$ is the greatest common
divisor of $\gamma$ and $\mu$, and $\chi(n)$ is an $f$th root of unity.

2) For $q=2$, $m>4$ (we do not take the case $m=4$ into account
here) every odd number $n$ has two indices $\alpha, \beta$, which are determined
with respect to the moduli $2, \frac{1}{2}\mu$, for which
\weq{7}{n \equiv (-1)^{\alpha}\,5^{\beta}\ (\mathrm{mod}\ m)}
holds. As characters one obtains, when $\varepsilon = \pm 1$ and $\Theta$ is set equal
to a $\frac{1}{2}\mu$th root of unity:
\weq{8}{\chi(n) = \varepsilon^{\alpha}\Theta^{\beta};}
$f$ is the least positive solution of the two congruences:
\[
\alpha f \equiv 0\ (\mathrm{mod}\ 2), \quad \beta f \equiv 0\ (\mathrm{mod}\ \tfrac{1}{2}\mu),
\]
and here too $\chi(n)$ is an $f$th root of unity.

For both cases, however, the following remark holds:

3) If $n$ belongs to the exponent $f$, then, if one lets $\chi$
run through the entire group of the $\mu$ characters, one obtains from
$\chi(n)$ every $f$th root of unity $e$ times.

For since the $\chi$ form an \gesperrt{Abel}ian group, one obtains
in the first place every $f$th root of unity that occurs at all
equally often, namely as often as the value 1, because, if $\chi_1(n)=\chi_2(n)$,
then $\chi_1\chi_2^{-1}(n) = \chi_0(n) = 1$. On the other hand one obtains
%CONT>>
